\subsection{LIE SUBGROUPS}

Let G be a Lie group and H a subgroup of G which is at the same time a submanifold of G. Then the mapping \((x,y)\mapsto xy^{-1}\) of \(H\times H\) into G is analytic and hence the mapping \((x,y)\mapsto xy^{-1}\) of \(H\times H\) into H is analytic. (Differentiable and Analytic Manifolds, R, 5.8.5). Thus H, with the group and manifold structures induced by those on G, is a Lie group.

DEFINITION 3. Let G be a Lie group. A subset H of G is called a Lie subgroup if H is a subgroup and a submanifold of G.

An open subgroup of G is a Lie subgroup of G. In particular, if G is a real or complex Lie group, its identity component is a Lie subgroup of G.

PROPOSITION 5. Let G be a Lie group and H a Lie subgroup of G.

\begin{enumerate}
\def\labelenumi{(\roman{enumi})}
\item
  H is closed in G.
\item
  The canonical injection of H into G is a Lie group morphism.
\item
  Let \(\mathbf{L}\) be a Lie group and \(f\) a mapping of \(\mathbf{L}\) into \(\mathbf{G}\) such that \(f(\mathbf{L}) \subset \mathbf{H}\) . For \(f\) to be a morphism of \(\mathbf{L}\) into \(\mathbf{H}\) , it is necessary and sufficient that \(f\) be a morphism of \(\mathbf{L}\) into \(\mathbf{G}\) .
\end{enumerate}

By Differentiable and Analytic Manifolds, R, 5.8.3, H is locally closed. Hence H is closed (General Topology, Chapter III, § 2, Proposition 4). Assertion (ii) is obvious. Assertion (iii) follows from Differentiable and Analytic Manifolds, R, 5.8.5.

PROPOSITION 6. Let G be a Lie group and H a subgroup of G. For H to be a Lie subgroup of G, it is necessary and sufficient that there exist a point \(h \in H\) and an open neighbourhood U of h in G such that \(H \cap U\) is a submanifold of G.

The condition is obviously necessary. Suppose that it holds. For all \(h' \in \mathbf{H}\) , the translation \(\gamma (h'h^{-1})\) is an automorphism of the manifold \(G\) and maps the submanifold \(H \cap U\) of \(U\) into the submanifold \((h'h^{-1}H) \cap (h'h^{-1}U)\) of \(h'h^{-1}U\) . As \(h'h^{-1}H = H\) and \(h'h^{-1}U\) is an open neighbourhood of \(h'\) in \(G\) , we see that every point of \(H\) has an open neighbourhood \(V\) such that \(V \cap H\) is a submanifold of \(G\) . Hence \(H\) is a submanifold of \(G\) .

Let \(G\) be a Lie group and \(H\) a Lie subgroup of \(G\) . If \(L\) is a Lie subgroup of \(H\) , \(L\) is a Lie subgroup of \(G\) by Differentiable and Analytic Manifolds, \(R\) , 5.8.6. Let \(M\) be a Lie subgroup of \(G\) such that \(M \subset H\) . Then \(M\) is a Lie subgroup of \(H\) , for the canonical injection of \(M\) into \(H\) is obviously an immersion.

Let k be a non-discrete closed subfield of K. A Lie k-subgroup of G is a Lie subgroup of the underlying Lie k-group of G.

Remark. If ``submanifold'' is replaced by ``quasi-submanifold'' in Definition 3, we obtain the definition of Lie quasi-subgroups of G. (For finite-dimensional G, the Lie quasi-subgroups are just the Lie subgroups.) Suppose that K is of characteristic 0. Proposition 5 remains valid with the same proof for Lie quasi-subgroups. Proposition 6 remains valid with the same proof, replacing ``Lie subgroup'' by ``Lie quasi-subgroup'' and ``submanifold'' by ``quasi-submanifold''.

